\chapter{Introduction}
\label{ch:introduction}

Understanding why some individuals develop a disease while others do not
is one of the central questions in modern medicine. For many common
diseases, including diabetes, heart disease, and psychiatric disorders,
the answer lies partly in genetics: individuals inherit different
combinations of genetic variants, and these variants collectively
influence their risk of disease. If genetic risk could be accurately
predicted from an individual's DNA, it would enable earlier
interventions, targeted screening, and personalised treatment strategies
\citep{Torkamani2018}.

Over the past two decades, genome-wide association studies (GWAS) have
identified thousands of genetic variants associated with complex traits
\citep{Visscher2017}. Each variant typically contributes only a small
amount to disease risk, but when aggregated across the genome, these
contributions can be substantial. The standard tool for this aggregation
is the polygenic risk score (PRS), which sums the estimated effects of
many variants weighted by an individual's genotype
\citep{Choi2020}. Bayesian approaches such as LDpred2 represent the
current state of the art, accounting for the correlation structure
between nearby variants and using prior assumptions about the genetic
architecture to produce more accurate predictions \citep{Prive2023}.

However, all existing PRS methods share a common limitation: they assume
that genetic variants contribute to disease risk in a purely additive
manner. There is growing evidence that nonlinear and context-dependent
interactions between variants play a role in disease risk
\citep{Cordell2009}, and capturing these interactions requires models
that go beyond the additive framework. Transformer models, originally
developed for natural language processing \citep{AttentionIsAllYouNeed},
offer a potential path forward. These models learn complex dependencies
within sequential data through self-attention, and genetic data has a
similar sequential structure: the dependencies between nearby variants,
known as linkage disequilibrium, are analogous to the dependencies
between words in a sentence \citep{SNPbag}.

Applying transformers to genetic data is, however, primarily a
computational challenge. A typical genome-wide dataset contains millions
of variants across hundreds of thousands of individuals, and the memory
requirements of transformer models scale with sequence length. Training
such models requires significant GPU infrastructure, careful memory
management, and practical decisions about batch sizes, sequence lengths,
and training schedules that are all constrained by the available
hardware. This thesis is therefore a computational statistics project in
equal measure to a modelling project: a substantial portion of the work
has been dedicated to implementing, profiling, and scaling the training
pipeline across multiple GPUs, to data handling and preprocessing at
scale, and to understanding the tradeoffs between model capacity, data
throughput, and wall-clock time. The modelling decisions described in
this thesis are inseparable from these engineering constraints.\\

The purpose of this thesis is to develop and evaluate SNPbag, a
transformer-based model for binary phenotype prediction from genotype
data. The evaluation involves comparing the predictive performance of
SNPbag against LDpred2-auto \citep{Prive2023} on simulated phenotypes
with known genetic architecture, using classification accuracy and AUC
on a shared held-out test set.